\section{Вычислительные эксперименты}\label{sec4}

В данной главе проводится эмпирическая проверка теоретических оценок главы~\ref{sec3}. Рассматриваются две модели существенно разного масштаба: полносвязная нейронная сеть для классификации рукописных цифр, в которой можно непосредственно вычислять все три критерия $\Delta_1$, $\Delta_2$ и $\Delta_2^{(D)}$ и сопоставлять их зависимости от объема выборки с теоретическими порядками $\mathcal O(k^{-1})$ и $\mathcal O(k^{-2})$ (раздел~\ref{subsec:exp-mlp}); и предобученная трансформерная языковая модель порядка $10^8$ параметров, на которой проверяются содержательные свойства подпространственного критерия $\Delta_2^{(D)}$: согласие с полнопространственным критерием при малых масштабах возмущения, локальная область применимости квадратичной аппроксимации и относительные затраты трех предложенных алгоритмов оценивания (раздел~\ref{subsec:exp-transformer}).

Программная реализация всех экспериментов выполнена на фреймворке PyTorch и доступна по адресу:
\begin{center}
    \href{https://github.com/intsystems/Kiselev-MS-Thesis/tree/main/code}{\texttt{https://github.com/intsystems/Kiselev-MS-Thesis/tree/main/code}}.
\end{center}

\subsection{Полносвязная нейронная сеть на MNIST}\label{subsec:exp-mlp}

\textbf{Постановка эксперимента.} Эксперименты проводятся на полносвязной нейронной сети для классификации рукописных цифр на датасете MNIST. Архитектура содержит один скрытый слой ширины $64$ с ReLU-активацией и выходной линейный классификатор. Размеры обучающей выборки рассматриваются на логарифмически распределенной сетке
\[
k \in \{100, 150, 200, 300, 400, 500, 700, 1000, 1500, 2000, 3000, 5000, 7000, 10000\}.
\]

Для каждого размера выборки $k$ модель обучается в течение $100$ эпох оптимизатором Adam с шагом обучения $10^{-3}$, после чего в полученной точке параметров вычисляются критерии $\Delta_1(k)$, $\Delta_2(k)$ и $\Delta_2^{(D)}(k)$ для $D \in \{5, 10, 20\}$. Параметры пробных распределений: радиус окрестности $R=0{,}5$ для одноточечного критерия (усреднение по 100 случайным направлениям), стандартное отклонение $\sigma=0{,}5$ для полнопространственного и подпространственных среднеквадратичных критериев ($200$ точек Монте-Карло). Для вычисления подпространственных критериев предварительно строится подпространство наибольшей кривизны через дефлированную степенную итерацию на произведениях матрицы Гессе на вектор. Эксперимент повторяется для трех независимых инициализаций (seed $42$, $43$, $44$), результаты усредняются по запускам.

\textbf{Скорость убывания критериев.} На рисунке~\ref{fig:landscape-convergence} представлена зависимость критериев $\Delta_1(k)$, $\Delta_2(k)$ и $\Delta_2^{(D)}(k)$ от размера выборки $k$ в логарифмических координатах. Пунктирные линии соответствуют теоретическим показателям $\propto k^{-1}$ и $\propto k^{-2}$.

\begin{figure}[ht]
    \centering
    \includegraphics[width=0.9\linewidth]{landscape_convergence.pdf}
    \caption{Зависимость критериев $\Delta_1(k)$, $\Delta_2(k)$ и $\Delta_2^{(D)}(k)$ от размера обучающей выборки $k$ в логарифмических координатах. Пунктирные линии соответствуют теоретическим порядкам $\propto k^{-1}$ и $\propto k^{-2}$}
    \label{fig:landscape-convergence}
\end{figure}

Для количественной проверки полученных оценок к каждой зависимости подгоняется степенной закон $\Delta(k) \sim k^{-\alpha}$. Оценки показателей и коэффициентов детерминации $R^2$ приведены в таблице~\ref{tab:landscape-exponents}.

\begin{table}[ht]
    \centering
    \caption{Оцененные показатели степенного закона $\Delta(k) \sim k^{-\alpha}$ для критериев на MLP/MNIST}
    \label{tab:landscape-exponents}
    \begin{tabular}{@{}lccc@{}}
        \toprule
        Критерий & Оценка $\alpha$ & $R^2$ & Теоретический порядок \\
        \midrule
        $\Delta_1$ & $2{,}33$ & $0{,}63$ & $\geqslant 1$ \\
        $\Delta_2$ & $1{,}94$ & $1{,}00$ & $\approx 2$ \\
        $\Delta_2^{(D=5)}$ & $1{,}80$ & $0{,}96$ & $\approx 2$ \\
        $\Delta_2^{(D=10)}$ & $1{,}83$ & $0{,}99$ & $\approx 2$ \\
        $\Delta_2^{(D=20)}$ & $1{,}84$ & $0{,}99$ & $\approx 2$ \\
        \bottomrule
    \end{tabular}
\end{table}

Полученные численные показатели согласуются с теоретическими порядками:

\begin{itemize}
    \item Среднеквадратичный критерий $\Delta_2$ убывает гладко и хорошо описывается степенным законом с показателем, близким к $2$, что в точности соответствует теоретической оценке~\eqref{eq:squared-criterion-general-theorem}.
    \item Подпространственные критерии $\Delta_2^{(D)}$ для $D = 5, 10, 20$ также убывают близко к квадратичному закону ($\alpha$ от $1{,}80$ до $1{,}84$), что подтверждает теоретическую оценку~\eqref{eq:squared-subspace-general-theorem}. При этом значения подпространственных критериев систематически меньше полнопространственного $\Delta_2$, что согласуется с эмпирически наблюдаемой концентрацией кривизны в подпространстве ведущих собственных направлений.
    \item Одноточечный критерий $\Delta_1$ демонстрирует существенно более высокую дисперсию, чем усредненные квадратичные критерии (что отражено в более низком $R^2 = 0{,}63$), однако и для него убывание заведомо не медленнее теоретической оценки $\mathcal O(k^{-1})$.
\end{itemize}

\textbf{Анизотропия ландшафта.} Для качественного понимания геометрии поверхности функции потерь и обоснованности перехода к подпространственному критерию строятся дополнительные визуализации. На рисунке~\ref{fig:hessian-spectrum} представлены ведущие собственные значения матрицы Гессе для модели, обученной на подвыборке размера $k = 1000$, в контрольных точках обучения на $10$, $20$, $50$ и $100$ эпохах. Полулогарифмическая шкала показывает быстрое убывание собственных значений по индексу, что характерно для нейронных сетей и подтверждает локализацию кривизны в подпространстве сравнительно небольшой размерности.

\begin{figure}[ht]
    \centering
    \includegraphics[width=0.8\linewidth]{hessian_spectrum.pdf}
    \caption{Собственные значения матрицы Гессе $\lambda_i$ (логарифмическая шкала) на разных эпохах обучения для подвыборки $k=1000$. Убывание по индексу собственного значения близко к экспоненциальному}
    \label{fig:hessian-spectrum}
\end{figure}

На рисунке~\ref{fig:loss-surface-2d} представлены двумерные сечения функции потерь обученной модели вдоль двух случайных ортогональных направлений (слева) и вдоль двух ведущих собственных векторов матрицы Гессе (справа). Сечение вдоль собственных векторов демонстрирует выраженную анизотропию: в этих направлениях изменение функции потерь происходит значительно быстрее, чем в произвольных ортогональных направлениях. Это подтверждает, что ограничение пробной области на подпространство наибольшей кривизны (раздел~\ref{subsec:subspace}) согласуется с эмпирической геометрией ландшафта.

\begin{figure}[ht]
    \centering
    \includegraphics[width=\linewidth]{loss_surface_2d.pdf}
    \caption{Двумерные сечения ландшафта функции потерь: случайные ортогональные направления (слева) и ведущие собственные векторы матрицы Гессе (справа), $k=1000$}
    \label{fig:loss-surface-2d}
\end{figure}

\subsection{Трансформерная языковая модель}\label{subsec:exp-transformer}

\textbf{Постановка эксперимента.} Для проверки полученных результатов в моделях большой размерности используется предобученная декодерная трансформерная языковая модель \texttt{nanochat} глубины $6$ ($\sim 10^8$ параметров) с роторными позиционными эмбеддингами, групповым вниманием, ReLU-активациями и RMSNorm без обучаемых параметров. Контрольная точка модели соответствует $3500$-му шагу обучения. Этот выбор обусловлен тем, что при $\sim 10^8$ параметров повторные произведения матрицы Гессе на вектор в точности float32 еще остаются вычислительно осуществимыми, при этом модель сохраняет нетривиальную структуру второй производной, характерную для глубоких языковых моделей. Все вычисления выполняются в float32 без автокаста; для устойчивости численных операций attention используется математический SDPA-ядро.

В этой постановке буква $k$ обозначает число обучающих последовательностей, определяющих функцию потерь $\mathcal L_k$. Если не оговорено иное, $\mathcal L_k$ и $\mathcal L_{k+1}$ строятся при $k=8$, а подпространство наибольшей кривизны в точке $\mathbf w_k^*$ строится дефлированной степенной итерацией на произведениях матрицы Гессе на вектор.

\textbf{Согласие подпространственного критерия с полнопространственным.} Центральный качественный результат теоремы~\ref{thm:squared-subspace-general} состоит в том, что переход к подпространству наибольшей кривизны сохраняет порядок убывания критерия $\mathcal O(k^{-2})$ и заменяет зависимость от полной размерности $N$ на зависимость от размерности подпространства $D$. Для трансформерной модели типичные значения $D$ при экспериментах не превышают нескольких десятков, тогда как $N \sim 10^8$, поэтому $D/N < 10^{-6}$.

На рисунке~\ref{fig:exp-subspace-vs-full} приведено отношение $\Delta_2^{(D)}(k)/\Delta_2(k)$ как функция размера выборки $k$ при нескольких значениях $D$ и $\sigma$. При $\sigma = 10^{-4}$ и $\sigma = 10^{-3}$ это отношение близко к единице во всем диапазоне $k$, т.\,е. подпространственный критерий воспроизводит полнопространственный сигнал стабилизации в пределах численной погрешности при ничтожно малой доле направлений. При $\sigma = 10^{-2}$ отношение существенно отклоняется от единицы и заметно зависит от $D$, что соответствует выходу за пределы локального режима применимости квадратичной аппроксимации.

\begin{figure}[ht]
    \centering
    \includegraphics[width=\linewidth]{delta_k_curriculum_d6_8seeds_gm_vs_full_ratio.pdf}
    \caption{Отношение подпространственного критерия к полнопространственному $\Delta_2^{(D)}/\Delta_2$ в зависимости от размера выборки $k$ для нескольких значений $D$ и $\sigma$. При $\sigma=10^{-4}, 10^{-3}$ отношение близко к единице во всем диапазоне $k$, при $\sigma = 10^{-2}$ "--- расходится с $\Delta_2$ из-за выхода из локального режима}
    \label{fig:exp-subspace-vs-full}
\end{figure}

\textbf{Область применимости квадратичной аппроксимации.} Алгоритмы~\eqref{eq:quad-subspace-mc} и~\eqref{eq:quad-closed-form-expanded} опираются на локальную квадратичную модель приращения функции потерь. Чтобы определить диапазон масштабов возмущения $\sigma$, в котором эта модель достаточно точна, в фиксированной точке $\mathbf w_k^*$ случайно генерируются изотропные возмущения $\bm \delta \sim \mathcal N(\mathbf 0, \sigma^2 \mathbf I_N)$, и истинное приращение $\mathcal L_k(\mathbf w_k^* + \bm \delta)-\mathcal L_k(\mathbf w_k^*)$ сопоставляется с его квадратичной аппроксимацией $\mathbf g^{(k)\top}\bm \delta + \tfrac12 \bm \delta^\top \mathbf H^{(k)} \bm \delta$.

\begin{figure}[ht]
    \centering
    \includegraphics[width=0.7\linewidth]{quadratic_sigma_loggrid_1e7_1e1_plot_abs.pdf}
    \caption{Относительная погрешность квадратичной аппроксимации приращения функции потерь в зависимости от масштаба возмущения $\sigma$ (\texttt{nanochat d6}, шаг $3500$). До $\sigma \approx 10^{-3}$ погрешность остается малой и почти постоянной; при больших $\sigma$ резко растет}
    \label{fig:exp-quadratic-validity}
\end{figure}

Результат, приведенный на рисунке~\ref{fig:exp-quadratic-validity}, демонстрирует выраженный локальный режим: относительная погрешность квадратичной аппроксимации остается низкой и почти постоянной до $\sigma \approx 10^{-3}$, после чего резко возрастает. Поэтому алгоритмы, основанные на квадратичном суррогате (разделы~\ref{subsec:algorithms}), применяются в дальнейшем только при $\sigma \leqslant 10^{-3}$.

\textbf{Сравнение алгоритмов оценивания.} Сравнение трех алгоритмов оценивания подпространственного критерия из раздела~\ref{subsec:algorithms} "--- прямой Монте-Карло, квадратичной Монте-Карло и оценки замкнутой формы через моменты квадратичной формы "--- приведено в таблице~\ref{tab:exp-runtimes}. Указаны средние времена выполнения каждой из трех стадий: построение подпространства (общее для всех алгоритмов), сборка коэффициентов $a_k$, $\mathbf c_k$, $\mathbf B_k$ (для алгоритмов, основанных на квадратичном суррогате) и собственно вычисление оценки.

\begin{table}[ht]
    \centering
    \caption{Время выполнения трех алгоритмов оценивания критерия $\Delta_2^{(D)}$ на модели \texttt{nanochat d6} (шаг $3500$, $D=45$, $\sigma=10^{-3}$), секунды (среднее $\pm$ выборочное стандартное отклонение по пяти запускам). Стадия I "--- построение подпространства (общая); стадия II "--- сборка градиента и сжатой матрицы Гессе (только для алгоритмов 2 и 3); стадия III "--- вычисление оценки}
    \label{tab:exp-runtimes}
    \begin{tabular}{@{}lccc@{}}
        \toprule
        Алгоритм & Стадия I (с) & Стадия II (с) & Стадия III (с) \\
        \midrule
        Прямая Монте-Карло & $17{,}53 \pm 0{,}53$ & "--- & $2{,}197 \pm 0{,}013$ \\
        Квадратичная Монте-Карло & $17{,}53 \pm 0{,}53$ & $1{,}777 \pm 0{,}014$ & $(3{,}68 \pm 0{,}05)\cdot 10^{-3}$ \\
        Замкнутая форма & $17{,}53 \pm 0{,}53$ & $1{,}777 \pm 0{,}014$ & $(1{,}22 \pm 0{,}03)\cdot 10^{-4}$ \\
        \bottomrule
    \end{tabular}
\end{table}

После общего шага построения подпространства оценка замкнутой формы оказывается приблизительно в $1{,}8\cdot 10^4$ раз быстрее прямой Монте-Карло оценки и приблизительно в $30$ раз быстрее квадратичной Монте-Карло. При этом, как показано на рисунке~\ref{fig:exp-quadratic-validity} и из согласия рисунка~\ref{fig:exp-subspace-vs-full}, в локальном режиме $\sigma \leqslant 10^{-3}$ обе оценки квадратичного суррогата дают результат, согласующийся с прямой Монте-Карло оценкой истинного критерия в пределах численной погрешности. Следовательно, в этом режиме оценка~\eqref{eq:quad-closed-form-expanded} замкнутой формы является не только корректной, но и существенно более экономичной альтернативой прямой Монте-Карло оценке.

\textbf{Основные результаты главы.} Проведенные вычислительные эксперименты подтверждают основные теоретические результаты главы~\ref{sec3}.

\begin{enumerate}
    \item На полносвязной нейронной сети, обученной на подвыборках разного размера датасета MNIST, среднеквадратичный критерий $\Delta_2$ и подпространственные критерии $\Delta_2^{(D)}$ для $D=5, 10, 20$ убывают с ростом $k$ по степенному закону с показателем, близким к теоретическому значению $\alpha = 2$ (см. таблицу~\ref{tab:landscape-exponents}). Спектр матрицы Гессе и двумерные сечения ландшафта подтверждают локализацию кривизны в подпространстве сравнительно небольшой размерности и анизотропию поверхности вдоль ведущих собственных направлений.
    \item На трансформерной языковой модели \texttt{nanochat d6} ($\sim 10^8$ параметров) подпространственный критерий $\Delta_2^{(D)}$ воспроизводит полнопространственный сигнал стабилизации в пределах численной погрешности при $D/N < 10^{-6}$ во всем локальном режиме $\sigma \leqslant 10^{-3}$, что подтверждает практическую целесообразность ограничения пробной области на подпространство наибольшей кривизны.
    \item Локальная квадратичная аппроксимация приращения функции потерь точна до $\sigma \approx 10^{-3}$, что задает область применимости алгоритмов, основанных на квадратичном суррогате (раздел~\ref{subsec:algorithms}).
    \item В этой области применимости оценка замкнутой формы через моменты квадратичной формы оказывается приблизительно в $1{,}8\cdot 10^4$ раз быстрее прямой Монте-Карло оценки без потери точности (таблица~\ref{tab:exp-runtimes}), подтверждая практическую ценность предложенной алгоритмической схемы для моделей большой размерности.
\end{enumerate}

Таким образом, эмпирическая проверка теоретических оценок главы~\ref{sec3} согласуется как с предсказанным порядком убывания критериев, так и с предсказанной заменой зависимости от размерности пространства параметров $N$ на зависимость от размерности подпространства $D$ при переходе к подпространственному критерию.
